\documentclass[10pt]{article}

\usepackage[utf8]{inputenc}

\usepackage[T1]{fontenc}

\usepackage{amsmath}

\usepackage{amsfonts}

\usepackage{amssymb}

\usepackage[version=4]{mhchem}

\usepackage{stmaryrd}



\begin{document}

ПЛОСКАЯ ЗАДАЧА т е о р и и у п р у г о с т иназвание типа задач, в к-рых картина изучаемого явления в упругой среде одинакова во всех плоскостях, параллельных нек-рой плоскости (напр., плоскости $O x_{1} x_{2}$ декартовой системы координат $\left.O x_{1} x_{2} x_{3}\right)$. Математич. теорией П. З. часто описываются и задачи, к-рые по содержанию имеют пространственный характер (напр., изгиб тонких пластинок).



По П. З. теории упругости успехи достигнуты главным образом благодаря пспользованию формул, выражающих искомые решения через аналитич. функции одного комплексного аргумента; впервые эти формулы были выведены в 1909 Г. К. Колосовым (см. [1]), а начиная с 20-х гг. 20 в. в работах Н. И. Мусхелишвили они нашли полное обоснование, и на их базе развиты методы решения широкого круга граничных (краевых) и контактных П. з. теории упругости. Полученные по П. з. теоретич. результаты используются при решении практич. задач.



Комплексное представление полей смещений и напряжений. Говорят, что упругая среда находится в состоянии плоской деформации, если существует такая декартова система координат $O x_{1} x_{2} x_{3}$, относительно к-рой компоненты вектора смещения имеют вид



\[

u_{\alpha}=u_{\alpha}\left(x_{1}, x_{2}, t\right), \alpha=1,2, u_{3}=0 \text {, }

\]



где $t$ - время. Компоненты тензора напряжений таковы



\[

X_{\alpha \beta}=\lambda \theta \delta_{\alpha \beta}+2 \mu e_{\alpha \beta}, X_{\alpha_{3}}=0, X_{33}=\lambda \theta,

\]



где $\lambda, \mu$ - постоянные Ламе, $\delta_{\alpha \beta}$ - символы Кронекера, $e_{\alpha \beta}$ - компоненты тензора деформации: $e_{\alpha \beta}=$ $=\partial_{\alpha} u_{\beta}+\partial_{\beta} u_{\alpha} ; \theta=e_{\alpha \alpha}=\partial_{\alpha} u_{\alpha}-$ дилатация $(\alpha, \beta=1,2$; наличие в выражении двух одинаковых индексов означает суммирование).



Плоская деформация возможна в упругой среде, к-рая заполняет цилиндр с образующими, перпендикулярными плоскости $O x_{1} x_{2}$, если при этом компоненты объемных сил имеют вид $X_{\alpha}=X_{\alpha}\left(x_{1}, x_{2}, t\right), X_{3}=0$, а боковые усилия не зависят от координаты $x_{3}$ и расположены на плоскостях, перпендикулярных оси цилиндра. Для реализации плоской деформации упругого цилиндра к его торцам необходимо приложить нормальные силы, равные $\pm \lambda \theta$.



При сделанных допущениях система уравнений динамики упругого тела относительно компонентов вектора смещений имеет вид



\[

\mu \Delta u_{\alpha}+(\lambda+\mu) \partial_{\alpha} \theta+X_{\alpha}=\rho \ddot{u}_{\alpha}, \alpha=1,2,

\]



где $\rho$ - плотность распределения масс, $\rho \ddot{u ̈}_{\alpha}-$ силы инерции, $\Delta$ - оператор Лапласа. Если воспользоваться операциями комплексного дифференцирования: $2 \partial_{\bar{z}}=\partial_{1}+i \partial_{2}, \quad 2 \partial_{z}=\partial_{1}-i \partial_{2}\left(\partial_{\alpha}=\partial / \partial x_{\alpha}\right)$, то при отсутствии сил инерции (статич. задача) эту систему можно записать в виде одного (комплексного) уравнения:



\[

(\lambda+3 \mu) \partial_{z \bar{z}}^{2} u+(\lambda+\mu) \partial_{\bar{z}}^{2} \bar{z} \bar{u}+X=0,

\]



где $\quad u=u_{1}+i u_{2}, X=2^{-1}\left(X_{1}+i X_{2}\right)$.



Пусть область $S$, занятая упругой средой, представляет связную часть плоскости $O x_{1} x_{2}$, ограниченную одним или несколькими контурами $L_{0}, L_{1}, \ldots, L_{m}$ без общих точек, $L=L_{0}+L_{1}+\ldots+L_{m}-$ граница области $S$; точка $z=0$ принадлежит $S$.



Искомое решение уравнения равновесия выражается по формуле $u=u_{0}+T X$, где $T X-$ нек-рое частное решение, к-рое можно выразить в виде



\[

\begin{gathered}

T X=-x(\mu \pi(1+x))^{-1} \iint X(\zeta) \ln |\zeta-z| d \zeta_{1} d \zeta_{2}+ \\

+(2 \mu \pi(1+x))^{-1} \iint \bar{X}(\zeta-z)(\bar{\zeta}-\bar{z})^{-1} d \zeta_{1} \dot{d} \zeta_{2},

\end{gathered}

\]



а $u_{0}$ - общее решение однородного уравнения $(X=0)$, к-рое выражается по формуле



\[

u_{0}=K(\varphi, \psi ; x)=x \varphi(z)-z \overline{\varphi^{\prime}(z)}-\psi(z)

\]



где $\varphi$ и $\psi$ - произвольные аналитич. функции от $z=x_{1}+i x_{2}$ в области $S(x=3-4 \sigma$, где $\sigma-$ постоянная Пуассона, $0<\sigma<0,5)$. Если $X$ - многочлен от $x, y$, то $T X$ можно выразить в явной форме.



Оператор $K(\varphi, \psi ; x)$ не изменится, если функции $\varphi$ и $\psi$ подчинить условию $\varphi(0)=0$ или $\psi(0)=0$. При выполнении одного из этих условпй всякому полю смецений $u=u_{1}+i u_{2}$, заданному в области $S$, соответствует вполне определенная пара аналитич. функций $\varphi, \psi$.



Если в предыдущих формулах постоянную $x$ заменить через $x^{*}=(3-\sigma) /(1+\sigma)$, то получается формула для поля смещения обобщенного плосконапряженного состояния.



Комплексная запись



$X_{\alpha \alpha}=2(\lambda+\mu)\left(\partial_{z} u+\partial_{-} u\right), X_{11}-X_{22}+2 i X_{12}=4 \mu \partial_{-} u$



компонентов тензора напряжений, в силу равенства



\[

u=K(\varphi, \psi ; x)+T x,

\]



принимает вид



\[

\begin{gathered}

X_{\alpha \alpha}=4 \operatorname{Re} \Phi(z)+T_{0} X, \\

\left.X_{11}-X_{22}+2 i X_{12}=2 \overline{(z} \Phi^{\prime}(z)+\Psi(z)\right)+T_{1} X, \\

\Phi=2 \mu \varphi^{\prime}, \Psi=2 \mu \Psi^{\prime}, \\

T_{0}=4(\lambda+\mu) \operatorname{Re} \partial_{z} T X, \quad T_{1} X=4 \mu \partial_{z} T X .

\end{gathered}

\]



Пусть упругая среда подвергается непрерывной деформации. Тогда можно считать, что компоненты тензора напряжений и смещений - непрерывные однозначные функции в области $S$; Ф и $\Psi$ голоморфны в $S$, причем $\Phi$ можно подчинить условию $\Phi^{\prime}(0)=\overline{\Phi^{\prime}(0)}$.



Если область $S$ ограничена и односвязна, а деформация непрерывна, то функции $\varphi$ и $\psi$ голоморфны в $S$. В случае конечной многосвязной области $\varphi$ и $\psi$ будут, вообще говоря, многозначными функциями определенного вида.



Основными задачами плоской теории упругости являются следующие задачи.



\begin{enumerate}

  \item П е р в а о сно в н а я 3 а д а ч а : определить упругое равновесие тела, когда на его границе заданы внепние силы.

\end{enumerate}



Сила напряжения $\left(X_{n}, Y_{n}\right)$, действующая на әлемент дуги $d s$ контура $L$ с нормалью $n$, может быть записана в комплексной форме:



\[

\left(X_{n}+i Y_{n}\right) d s=-2 i \mu d\left(\varphi(z)+z \overline{\varphi^{\prime}(z)}+\overline{\psi(z))},\right.

\]



и краевые условия первой задачи имеют вид



\[

\varphi(t)+t \overline{\varphi^{\prime}(t)}+\overline{\psi(t)}=f(t)+c(t), t \in L,

\]



где



\[

f(t)=i(2 \mu)^{-1} \int_{L}\left(X_{n}+i Y_{n}\right) d s,

\]



причем дуга $s$ отсчитывается на каждом $L_{k}$ от нек-рой фиксированной точки $z_{k} \in L_{k}$ в положительном направлении; $c(t)=c_{k}=$ const на $L_{k}$. Всегда можно считать $c_{0}=0$, остальные постоянные $c_{k}$ определяются в ходе решения задачи. Если $m=0$, то искомые функции $\varphi$ и $\psi$ голоморфны в $S$. Тогда равенства $\varphi(0)=0, \operatorname{Im} \varphi^{\prime}(0)=$ $=0$ обеспечивают единственность решения задачи (1), а необходимые и достаточные условия существования решения



\[

\begin{aligned}

\int_{L}\left(X_{n}+i Y_{n}\right) d s & =0,2 \mu \int_{L}\left(x_{1} Y_{n}-x_{2} X_{n}\right) d s= \\

& =\operatorname{Re} \int_{L} f \overline{d t}=0

\end{aligned}

\]





\end{document}